\chapter{Conclusion et Perspectives}

Le projet \textbf{\projectName} a atteint ses objectifs initiaux en livrant une solution de gestion de caisse sécurisée, moderne et centralisée. En remplaçant les processus manuels par un système numérique rigoureux, l'application réduit considérablement les risques d'erreurs et de fraudes.

L'architecture choisie (Next.js + SQLite) s'est avérée être un excellent compromis entre rapidité de développement, performance et maintenabilité. La séparation claire entre le Front-office (Caissier) et le Back-office (Admin) garantit une expérience utilisateur adaptée à chaque profil.

\section*{Perspectives d'Évolution}
Bien que le système soit pleinement opérationnel, plusieurs pistes d'amélioration sont envisagées pour les versions futures :

\begin{itemize}
    \item \textbf{Intégration TPE (Terminaux de Paiement) :} Connecter directement l'application aux terminaux bancaires pour automatiser la saisie des montants payés par carte, éliminant une autre source d'erreur de saisie.
    \item \textbf{Application Mobile :} Développer une version mobile (PWA ou React Native) pour permettre aux administrateurs de surveiller les agences depuis leur smartphone en déplacement.
    \item \textbf{Analyse Prédictive :} Utiliser l'historique des données pour prédire les besoins en fond de caisse par jour et par agence, optimisant ainsi la trésorerie immobilisée.
\end{itemize}

Ce projet constitue donc une fondation solide pour la transformation digitale de la gestion financière du réseau d'agences.
